\section{数据与文件格式}
\label{sec:formats}

\subsection{用户输入：YAML 参数化模型}

顶层字段与约束：
\begin{itemize}
  \item 必填：\texttt{oasis\_version}、\texttt{theme}，以及至少一个模型定义块（\texttt{source} 必选；
        \texttt{background}、\texttt{transmission} 可选）；编码 UTF-8；
  \item 物理量统一以 \texttt{\{value, unit\}} 表达，单位兼容天文标准单位制；
  \item \texttt{extends} 支持引用基础模型或上级主题；参数扫描支持 \texttt{linspace} / \texttt{logspace} / \texttt{list}；
  \item 物理语义校验（量纲一致性、非物理取值）由模型段执行，不前置在文件解析层。
\end{itemize}

示例（节选）：
\begin{verbatim}
oasis_version: "1.0"
theme: "stellar"
source:
  name: "sun_like_1"
  extends: "stellar_base"
  coordinates:
    ra: {value: 83.633, unit: "deg"}
    dec: {value: 22.014, unit: "deg"}
    distance: {value: 10, unit: "pc"}
  proper_motion:
    pm_ra_cosdec: {value: 34.5, unit: "mas/yr"}
    pm_dec: {value: -12.3, unit: "mas/yr"}
  spectral_model:
    type: "blackbody"
    parameters:
      temperature: {value: 5778, unit: "K"}
transmission:
  name: "ism_extinction"
  optical_depth: {value: 0.3, unit: ""}
  spectral_model:
    type: "power_law"
    parameters:
      alpha: {value: -1.5, unit: ""}
      reference_wavelength: {value: 5500, unit: "angstrom"}
\end{verbatim}

\subsection{正演输出包（目标层 $\rightarrow$ observables）}

\begin{itemize}
  \item \textbf{主 FITS}：\texttt{TargetHistogram} 序列化（二维/三维/四维数组），WCS 完备自洽，
        光变曲线含时间轴描述；
  \item \textbf{必填头关键字}：\texttt{OASIS\_VER}、\texttt{MODEL\_ID}、\texttt{TIMESTAMP}、
        \texttt{TARGET\_CLASS}、\texttt{OASIS\_OBS\_INTF\_VER}，以及 \texttt{SRC\_MODEL} /
        \texttt{BKG\_MODEL} / \texttt{TRANS\_MODEL} / \texttt{SPATIAL\_TYP}（无对应内容时为空字符串）；
  \item \textbf{辅助文件}：\texttt{params.json}（参数侧写）、\texttt{config.yaml}（观测配置摘要）、
        \texttt{processing\_log.txt}（计算过程日志）；CSV/TSV 处理指令可选；
  \item \textbf{哨兵}：全部写入完成后生成 \texttt{.forward\_ready}。
\end{itemize}

\subsection{反演输入包（observables $\rightarrow$ 目标层）}

\begin{itemize}
  \item \textbf{主 FITS}：科学数据产品，结构与正演输出一致（\texttt{TargetHistogram} 存储形态）；
  \item \textbf{头关键字}：\texttt{TASK\_ID}、\texttt{INSTRUMENT}、\texttt{CALIB\_VER}、\texttt{DATA\_QUAL}
        （good / warning / bad）、\texttt{OASIS\_OBS\_INTF\_VER}；
  \item \textbf{辅助文件}：\texttt{noise\_model.json}（噪声模型）、\texttt{observation\_config.yaml}
        （观测配置）、\texttt{instrument\_response.json}（仪器响应，可选）；
  \item \textbf{哨兵}：\texttt{.inverse\_ready}；
  \item \textbf{异常路径}：\texttt{.error} 文件承载 JSON 载荷（\texttt{error\_code} / \texttt{english\_message} /
        \texttt{timestamp}），错误码统一为 \texttt{E\_\{DOMAIN\}\_\{CODE\}} 格式（如 \texttt{E\_OBS\_001}）。
\end{itemize}

\subsection{哨兵协议}

\begin{center}
\begin{tabular}{lll}
\toprule
哨兵文件 & 放置方 & 含义 \\
\midrule
\texttt{.forward\_ready} & 目标层 & 正演产物就绪，observables 层可消费 \\
\texttt{.inverse\_ready} & observables 层 & 反演输入数据就绪 \\
\texttt{.error} & 任意方 & 处理失败，阻止后续自动流转 \\
\bottomrule
\end{tabular}
\end{center}

\subsection{与 OIFITS 2 的关系}

目标层与 observables 层之间交换的是目标区域的亮度分布产品（图像/光谱/光变），以 FITS 主数据 +
JSON/YAML 辅助文件表达，\emph{不使用 OIFITS 表结构}。OIFITS 2 是干涉观测量（可见度、闭合相位等）的交换标准，
适用于观测量层与载荷层之间及对外发布环节；当观测流程需要将目标模型转换为干涉观测量时，
由观测量层负责 OIFITS 封装，目标层不直接产出 OIFITS。
